\chapter{LIMITATIONS AND FUTURE ENHANCEMENTS}

\section{Limitations}
The current Smart Gym Assistant demonstrates exercise and posture recognition using
wearable inertial sensors, but several limitations remain:
\begin{itemize}
\item The implemented classifier supports only bicep curls and squats with the
      available labeled classes.
\item Recognition depends on consistent sensor placement and orientation on the
      user's body.
\item The training data may not represent every user's movement style, body type,
      speed, or range of motion.
\item Bad and good posture classes can overlap, particularly for squats, which
      reduces recall for incorrect posture detection.
\item The ESP32 uses a compact prototype set to limit memory and computation, so
      it may not capture every variation in the motion data.
\item Live dashboard monitoring depends on Wi-Fi availability and communication
      with the FastAPI server.
\end{itemize}

\section{Future Improvements}
The following improvements can extend the accuracy, usability, and scalability of
future versions of the system.

\subsection{Support for Additional Exercises}
Additional exercises such as push-ups, lunges, shoulder presses, and jumping jacks
can be added by collecting representative labeled data for each exercise and posture
condition. The training pipeline and embedded label definitions can then be extended
without changing the basic sensing architecture.

\subsection{Training with More Users and Motion Styles}
Future datasets should include multiple users, sensor placements, movement speeds,
and body-motion variations. This would improve generalization and reduce the risk
that the classifier performs well only for the user who collected the original data.

\subsection{Improved Posture Detection}
More carefully defined posture boundaries and additional bad-posture examples could
improve recall for safety-related classes such as \texttt{squat\_bad}. Future work
could also combine temporal features, movement ranges, and exercise-specific rules
with the KNN prediction.

\subsection{Optimized Embedded Model}
The model can be regenerated using the best validated prototype count and value of
$k$, provided the resulting memory and processing requirements remain suitable for
the ESP32. Quantization, dimensionality reduction, or a compact alternative model
could further reduce inference cost.

\subsection{Offline Data Buffering and Reliable Communication}
The ESP32 could temporarily buffer records when Wi-Fi is unavailable and upload
those records after reconnection. Retry handling, request status reporting, and
server-side validation could make collection and live monitoring more reliable.

\subsection{User Interface Enhancements}
The dashboard could provide session history, exercise summaries, repetition trends,
class-wise confidence, and calibration guidance. The OLED interface could also show
clearer status messages for sensor placement, Wi-Fi state, and posture feedback.
